\ifdefined\gkver\else\def\gkver{student}\fi
\documentclass[\gkver]{gaokaozhenti}
\begin{document}

\chapter{2025年广东}

\begin{enumerate}
\item
%% number: 1
%% paperName: 2025年广东高考第 1 题 4 分
%% typeId: 1
%% type: 单选题
%% chapter: 机械波
%% point: 多普勒效应
%% method: 无
%% score: 4
%% degree: 850
%% duplicateId: 0
%% body:
关于受迫振动和多普勒效应，下列说法正确的是\xzanswer{C}

\fourchoices[answer=C]
{系统的固有频率与驱动力频率有关}
{只要驱动力足够大，共振就能发生}
{应用多普勒效应可以测量车辆的速度}
{观察者与波源相互远离时，接收到的波的频率比波源的频率大}

%% answer: C
%% memo:
\memoanswer{
A．系统的固有频率只与系统本身有关，与驱动力频率无关，A 错误；

B．只有驱动力频率与系统固有频率相同时，共振才能发生，B 错误；

CD．根据多普勒效应可知观察者与波源相互远离时，接收到的波的频率比波源的频率小，观察者与波源相互靠近时，接收到的波的频率比波源的频率大，所以应用多普勒效应可以测量车辆的速度，C 正确，D 错误。

故选 C。
}

\item
%% number: 2
%% paperName: 2025年广东高考第 2 题 4 分
%% typeId: 1
%% type: 单选题
%% chapter: 电磁感应
%% point: 远距离输电
%% method: 无
%% score: 4
%% degree: 650
%% duplicateId: 0
%% body:
如图所示。某光伏电站输出功率 $1000 \UkW$、电压 $400 \UV$ 的交流电，经理想变压器升压至 $10 \UkV$ 后，通过输电线输送到变电站，输电线的等效电阻 $R$ 为 $5 \UO$。下列说法正确的是\xzanswer{B}

\onepicture[w=5.60cm]{figs/02.png}

\fourchoices[answer=B]
{变压器原、副线圈匝数比为 $1:100$}
{输电线上由 $R$ 造成的电压损失为 $500 \UV$}
{变压器原线圈中的电流为 $100 \UA$}
{变压器原、副线圈中电流的频率不同}

%% answer: B
%% memo:
\memoanswer{
A．根据理想变压器原副线圈电压比等于匝数比可得 $\frac{n_{1}}{n_{2}} = \frac{U_{1}}{U_{2}} = \frac{400}{10000} = \frac{1}{25}$，A 错误；

B．原副线圈两端的功率相等，流过副线圈的电流 $I_{2}=\frac{P}{U_{2}}=\frac{1000000}{10000}\UA=100\UA$

输电线上由 $R$ 造成的电压损失为 $\Delta U = I_{2}R = 100 \times 5\UV = 500\UV$，B 正确；

C．变压器原线圈中的电流为 $I_{1}=\frac{P}{U_{1}}=\frac{1000000}{400}\UA=2500\UA$，C 错误；

D．变压器不改变交变电流的频率，变压器原、副线圈中电流的频率相同，D 错误。

故选 B。
}

\item
%% number: 3
%% paperName: 2025年广东高考第 3 题 4 分
%% typeId: 1
%% type: 单选题
%% chapter: 光学
%% point: 光电效应
%% method: 无
%% score: 4
%% degree: 650
%% duplicateId: 0
%% body:
有甲、乙两种金属，甲的逸出功小于乙的逸出功。使用某频率的光分别照射这两种金属，只有甲发射光电子，其最大初动能为 $E_{k}$，下列说法正确的是\xzanswer{B}

\fourchoices[answer=B]
{使用频率更小的光，可能使乙也发射光电子}
{使用频率更小的光，若仍能使甲发射光电子，则其最大初动能小于 $E_{k}$}
{频率不变，减弱光强，可能使乙也发射光电子}
{频率不变，减弱光强，若仍能使甲发射光电子，则其最大初动能小于 $E_{k}$}

%% answer: B
%% memo:
\memoanswer{
A．某频率的光不能使乙金属发生光电效应，说明此光的频率小于乙金属的截止频率，则换用频率更小的光不能发生光电效应，A 错误；

B．由光电效应方程 $E_{k} = h\nu - W_{0}$ 可知频率越大最大初动能越大，换用频率更小的光最大初动能小于 $E_{k}$，B 正确；

C．频率不变则小于乙金属的截止频率，不会发生光电效应，C 错误；

D．由 $E_{k} = h\nu - W_{0}$ 可知频率不变最大初动能不变，D 错误。

故选 B。
}

\item
%% number: 4
%% paperName: 2025年广东高考第 4 题 4 分
%% typeId: 1
%% type: 单选题
%% chapter: 光学
%% point: 光的折射
%% method: 无
%% score: 4
%% degree: 650
%% duplicateId: 0
%% body:
如图为测量某种玻璃折射率的光路图。某单色光从空气垂直射入顶角为 $\alpha$ 的玻璃棱镜，出射光相对于入射光的偏转角为 $\beta$，该折射率为\xzanswer{A}

\onepicture[w=3.40cm]{figs/04.png}

\fourchoices[answer=A]
{$\frac{\sin(\alpha + \beta)}{\sin\alpha}$}
{$\frac{\sin(\alpha + \beta)}{\sin\beta}$}
{$\frac{\sin\alpha}{\sin\beta}$}
{$\frac{\sin\beta}{\sin\alpha}$}

%% answer: A
%% memo:
\memoanswer{
光路图如图所示。

\onepicture[w=3.40cm]{figs/04_an.png}

则有折射定律可得 $n = \frac{\sin(\alpha + \beta)}{\sin\alpha}$

故选 A。
}

\item
%% number: 5
%% paperName: 2025年广东高考第 5 题 4 分
%% typeId: 1
%% type: 单选题
%% chapter: 曲线运动
%% point: 开普勒定律
%% method: 无
%% score: 4
%% degree: 550
%% duplicateId: 0
%% body:
一颗绕太阳运行的小行星，其轨道近日点和远日点到太阳的距离分别约为地球到太阳距离的 5 倍和 7 倍。关于该小行星，下列说法正确的是\xzanswer{D}

\fourchoices[answer=D]
{公转周期约为 6 年}
{从远日点到近日点所受太阳引力大小逐渐减小}
{从远日点到近日点线速度大小逐渐减小}
{在近日点加速度大小约为地球公转加速度的 $\frac{1}{25}$}

%% answer: D
%% memo:
\memoanswer{
A．根据题意，设地球与太阳间距离为 $R$，则小行星公转轨道的半长轴为 $a=\frac{5R+7R}{2}=6R$

由开普勒第三定律有 $\frac{(6R)^{3}}{T_{2}^{2}} = \frac{R^{3}}{T_{\text{地}}^{2}}$

解得 $T = \sqrt{6^{3}T_{\text{地}}^{2}} = 6\sqrt{6}$ 年

故 A 错误；

B．从远日点到近日点，小行星与太阳间距离减小，由万有引力定律 $F = \frac{Gm_{1}m_{2}}{r^{2}}$ 可知，小行星受太阳引力增大，故 B 错误；

C．由开普勒第二定律可知，从远日点到近日点，小行星线速度逐渐增大，故 C 错误。

D．由牛顿第二定律有 $\frac{GMm}{r^{2}} = ma$

解得 $a = \frac{GM}{r^{2}}$

可知 $\frac{a_{\text{行}}}{a_{\text{地}}} = \frac{R^{2}}{(5R)^{2}} = \frac{1}{25}$

即小行星在近日点的加速度是地球公转加速度的 $\frac{1}{25}$，故 D 正确；

故选 D。
}

\item
%% number: 6
%% paperName: 2025年广东高考第 6 题 4 分
%% typeId: 1
%% type: 单选题
%% chapter: 磁场
%% point: 质谱仪回旋加速器
%% method: 无
%% score: 4
%% degree: 450
%% duplicateId: 0
%% body:
某同步加速器简化模型如图所示，其中仅直通道 PQ 内有加速电场，三段圆弧内均有可调的匀强偏转磁场 $B$。带电荷量为 $-q$、质量为 $m$ 的离子以初速度 $v_{0}$ 从 P 处进入加速电场后，沿顺时针方向在加速器内循环加速。已知加速电压为 $U$，磁场区域中离子的偏转半径均为 $R$。忽略离子重力和相对论效应，下列说法正确的是\xzanswer{A}

\onepicture[w=3.50cm]{figs/06.png}

\fourchoices[answer=A]
{偏转磁场的方向垂直纸面向里}
{第 1 次加速后，离子的动能增加了 $2qU$}
{第 $k$ 次加速后，离子的速度大小变为 $\frac{\sqrt{m^{2}v_{0}^{2} + kqUm}}{m}$}
{第 $k$ 次加速后，偏转磁场的磁感应强度大小应为 $\frac{\sqrt{m^{2}v_{0}^{2} - 2kqUm}}{qR}$}

%% answer: A
%% memo:
\memoanswer{
A．直线通道 PQ 有电势差为 $U$ 的加速电场，粒子带负电，粒子沿顺时针方向运动，由左手定则可知，偏转磁场的磁感应强度方向垂直纸面向里，故 A 正确；

BC．根据题意，由动能定理可知，加速一次后，带电粒子的动能增量为 $qU$，由于洛伦兹力不做功，则加速 $k$ 次后，带电粒子的动能增量为 $kqU$，加速 $k$ 次后，由动能定理有

$kqU = \frac{1}{2}mv^{2} - \frac{1}{2}mv_{0}^{2}$

解得 $v = \sqrt{v_{0}^{2} + \frac{2kqU}{m}} = \frac{\sqrt{m^{2}v_{0}^{2} + 2kqUm}}{m}$

故 BC 错误；

D．粒子在偏转磁场中运动的半径为 $R$，则有 $qvB = m\frac{v^{2}}{R}$

联立解得 $B=\frac{mv}{qR}=\frac{m}{qR}\sqrt{v_{0}^{2}+\frac{2kqU}{m}}=\frac{\sqrt{m^{2}v_{0}^{2}+2kmqU}}{qR}$

故 D 错误。

故选 A。
}

\item
%% number: 7
%% paperName: 2025年广东高考第 7 题 4 分
%% typeId: 1
%% type: 单选题
%% chapter: 运动和力
%% point: 动量和能量
%% method: 无
%% score: 4
%% degree: 350
%% duplicateId: 0
%% body:
如图所示，光滑水平面上，小球 M、N 分别在水平恒力 $F_{1}$ 和 $F_{2}$ 作用下，由静止开始沿同一直线相向运动，在 $t_{1}$ 时刻发生正碰后各自反向运动。已知 $F_{1}$ 和 $F_{2}$ 始终大小相等，方向相反。从开始运动到碰撞后第 1 次速度减为 0 的过程中，两小球速度 $v$ 随时间 $t$ 变化的图像，可能正确的是\xzanswer{A}

\onepicture[w=5.60cm]{\includesvg{figs/07a}}

\fourchoices[answer=A, ispicture=true, h=2.20cm]
{figs/07c.png}
{figs/07d.png}
{figs/07e.png}
{figs/07f.png}

%% answer: A
%% memo:
\memoanswer{
根据牛顿第二定律 $a = \frac{F}{m}$，两物体受外力 $F$ 大小相等，由图像的斜率等于加速度可知 M、N 的加速度大小之比为 $4:6 = 2:3$，可知 M、N 的质量之比为 $6:4 = 3:2$；设分别为 $3m$ 和 $2m$；由图像可设 MN 碰前的速度分别为 $4v$ 和 $6v$，则因 MN 系统受合外力为零，向右为正方向，则系统动量守恒，则由动量守恒定律

$3m \cdot 4v - 2m \cdot 6v = 3mv_{1} + 2mv_{2}$

若系统为弹性碰撞，则能量关系可知

$\frac{1}{2} \cdot 3m(4v)^{2} + \frac{1}{2} \cdot 2m(6v)^{2} = \frac{1}{2} \cdot 3mv_{1}^{2} + \frac{1}{2} \cdot 2mv_{2}^{2}$

解得 $v_{1} = -4v$、$v_{2} = 6v$

因 M、N 的加速度大小之比仍为 $2:3$，则停止运动的时间之比为 $1:1$，即两物体一起停止，则 BD 是错误的；

若不是弹性碰撞，则 $3m \cdot 4v - 2m \cdot 6v = 3mv_{1} + 2mv_{2}$

可知碰后速度大小之比为 $v_{1}:v_{2} = 2:3$

若假设 $v_{1} = 2v$，则 $v_{2} = 3v$，此时满足

$\frac{1}{2} \cdot 3m(4v)^{2} + \frac{1}{2} \cdot 2m(6v)^{2} > \frac{1}{2} \cdot 3mv_{1}^{2} + \frac{1}{2} \cdot 2mv_{2}^{2}$

则假设成立，因 M、N 的加速度大小之比仍为 $2:3$，则停止运动的时间之比为 $1:1$，对 M 来说碰撞前后的速度之比为 $4v:2v = 2:1$

可知碰撞前后运动时间之比为 $2:1$，可知 A 正确，C 错误。

故选 A。
}

\item
%% number: 8
%% paperName: 2025年广东高考第 8 题 6 分
%% typeId: 2
%% type: 多选题
%% chapter: 曲线运动
%% point: 圆周运动的应用
%% method: 无
%% score: 6
%% degree: 550
%% duplicateId: 0
%% body:
将可视为质点的小球沿光滑冰坑内壁推出，使小球在水平面内做匀速圆周运动，如图所示。已知圆周运动半径 $R$ 为 $0.4 \Um$，小球所在位置处的切面与水平面夹角 $\theta$ 为 $45^\circ$，小球质量为 $0.1 \Ukg$，重力加速度 $g$ 取 $10 \Umsq$。关于该小球，下列说法正确的有\xzanswer{AC}

\onepicture[w=6.00cm]{figs/08.png}

\fourchoices[answer=AC]
{角速度为 $5 \Urads$}
{线速度大小为 $4 \Ums$}
{向心加速度大小为 $10 \Umsq$}
{所受支持力大小为 $1 \UN$}

%% answer: AC
%% memo:
\memoanswer{
A．对小球受力分析可知 $F_{\text{向}}=mg\tan 45^\circ=m\omega^{2}R$

解得 $\omega = 5 \Urads$

故 A 正确；

B．线速度大小为 $v = \omega R = 2 \Ums$

故 B 错误；

C．向心加速度大小为 $a_{n} = \omega^{2}R = 10 \Umsq$

故 C 正确；

D．所受支持力大小为 $N=\frac{mg}{\cos 45^\circ}=\sqrt{2} \UN$

故 D 错误。

故选 AC。
}

\item
%% number: 9
%% paperName: 2025年广东高考第 9 题 6 分
%% typeId: 2
%% type: 多选题
%% chapter: 电磁感应
%% point: 切割磁感线电动势
%% method: 无
%% score: 6
%% degree: 550
%% duplicateId: 0
%% body:
如图是一种精确测量质量的装置原理示意图，竖直平面内，质量恒为 $M$ 的称重框架由托盘和矩形线圈组成。线圈的一边始终处于垂直线圈平面的匀强磁场中，磁感应强度不变。测量分两个步骤，步骤①：托盘内放置待测物块，其质量用 $m$ 表示，线圈中通大小为 $I$ 的电流，使称重框架受力平衡；步骤②：线圈处于断开状态，取下物块，保持线圈不动，磁场以速率 $v$ 匀速向下运动，测得线圈中感应电动势为 $E$。利用上述测量结果可得出 $m$ 的值，重力加速度为 $g$。下列说法正确的有\xzanswer{BD}

\onepicture[w=6.50cm]{figs/09.png}

\fourchoices[answer=BD]
{线圈电阻为 $\frac{E}{I}$}
{$I$ 越大，表明 $m$ 越大}
{$v$ 越大，则 $E$ 越小}
{$m = \frac{EI}{vg} - M$}

%% answer: BD
%% memo:
\memoanswer{
A．根据题意电动势 $E$ 是线圈断开时切割磁感线产生的感应电动势，$I$ 为线圈闭合时通入的电流，故 $\frac{E}{I}$ 不是线圈的电阻；

故 A 错误；

B．根据平衡条件有 $(M + m)g = BIL$ ①

故可知 $I$ 越大，$m$ 越大；

故 B 正确；

C．根据公式有 $E = BLv$ ②

故可知 $v$ 越大，$E$ 越大；

故 C 错误；

D．联立①②可得 $m = \frac{EI}{vg} - M$

故 D 正确。

故选 BD。
}

\item
%% number: 10
%% paperName: 2025年广东高考第 10 题 6 分
%% typeId: 2
%% type: 多选题
%% chapter: 力物体的平衡
%% point: 动态平衡
%% method: 无
%% score: 6
%% degree: 350
%% duplicateId: 0
%% body:
如图所示，无人机在空中作业时，受到一个方向不变、大小随时间变化的拉力。无人机经飞控系统实时调控，在拉力、空气作用力和重力作用下沿水平方向做匀速直线运动。已知拉力与水平面成 $30^\circ$ 角，其大小 $F$ 随时间 $t$ 的变化关系为 $F = F_{0} - kt$（$F \neq 0$，$F_{0}$、$k$ 均为大于 0 的常量），无人机的质量为 $m$，重力加速度为 $g$。关于该无人机在 0 到 $T$ 时间段内（$T$ 是满足 $F > 0$ 的任一时刻），下列说法正确的有\xzanswer{AB}

\onepicture[w=5.00cm]{figs/10.png}

\fourchoices[answer=AB]
{受到空气作用力的方向会变化}
{受到拉力的冲量大小为（$F_{0} - \frac{1}{2}kT$）$T$}
{受到重力和拉力的合力的冲量大小为 $mgT$ + （$F_{0} - \frac{1}{2}kT$）$T$}
{$T$ 时刻受到空气作用力的大小为 $\sqrt{\frac{3}{4}(F_{0} - kT)^{2} + (mg - \frac{F_{0} - kT}{2})^{2}}$}

%% answer: AB
%% memo:
\memoanswer{
AD．无人机经飞控系统实时调控，在拉力、空气作用力和重力作用下沿水平方向做匀速直线运动，则无人机受到空气作用力与重力和拉力的合力等大反向，随着 $F$ 的减小重力和拉力的合力如图。

\onepicture[w=3.00cm]{figs/10_an.png}

可知无人机受到空气作用力的大小和方向均会改变，在 $T$ 时刻有

$\cos 120^\circ=\frac{(mg)^{2}+F^{2}-F_{\text{空}}^{2}}{2mgF}$，$F=F_{0}-kT$

解得 $F_{\text{空}}=\sqrt{(F_{0}-kT)^{2}+(mg)^{2}+mg(F_{0}-kT)}$

故 A 正确、D 错误；

B．由于拉力 $F$ 随时间 $t$ 均匀变化，则无人机在 0 到 $T$ 时间段内受到拉力的冲量大小为 $F-t$ 图像与坐标轴围成的面积为（$F_{0} - \frac{1}{2}kT$）$T$，故 B 正确；

C．将拉力分解为水平和竖直方向，则无人机受重力和拉力的合力在水平方向有

$F_{x}=\frac{\sqrt{3}}{2}(F_{0}-kt)$

无人机受重力和拉力的合力在竖直方向有 $F_{y}=\frac{1}{2}(F_{0}-kt)+mg$

0 到 $T$ 时间段内无人机受重力和拉力的合力在水平方向的冲量为 $I_{x}=\frac{\sqrt{3}}{2}(F_{0}T-\frac{1}{2}kT^{2})$

0 到 $T$ 时间段内无人机受重力和拉力的合力在竖直方向的冲量为 $I_{y}=\frac{1}{2}F_{0}T-\frac{1}{4}kT^{2}+mgT$

则 0 到 $T$ 时间段内无人机受到重力和拉力的合力的冲量大小为

$I=\sqrt{I_{x}^{2}+I_{y}^{2}}=\sqrt{(F_{0}T-\frac{1}{2}kT^{2}+\frac{1}{2}mgT)^{2}+\frac{3}{4}m^{2}g^{2}T^{2}}$

故 C 错误。

故选 AB。
}

\item
%% number: 11
%% paperName: 2025年广东高考第 11 题 8 分
%% typeId: 4
%% type: 实验题
%% chapter: 运动和力
%% point: 动量守恒定律
%% method: 无
%% score: 8
%% degree: 550
%% duplicateId: 0
%% body:
请完成下列实验操作和计算。

\begin{enumerate}
\item
在“长度的测量及其测量工具的选用”实验中，用螺旋测微器测量小球的直径，示数如图~\ref{2025广东11a}所示，读数 \tkanswer{8.260} $\Umm$。
\item
实验小组利用小车碰撞实验测量吸能材料的性能，装置如图~\ref{2025广东11b}所示，图中轨道由轨道甲和乙平滑拼接而成，且轨道乙倾角较大。
\begin{steps}
\item
选取相同的两辆小车，分别安装宽度为 $1.00 \Ucm$ 的遮光条。
\item
轨道调节。
调节螺母使轨道甲、乙连接处适当升高。将小车在轨道乙上释放，若测得小车通过光电门 A 和 B 的 \tkanswer{时间相等}。证明已平衡小车在轨道甲上所受摩擦力及其他阻力。
\item
碰撞测试。
先将小车 1 静置于光电门 A 和 B 中间，再将小车 2 在 M 点由静止释放，测得小车 2 通过光电门 A 的时间为 $t_{2}$，碰撞后小车 1 通过光电门 B 的时间为 $t_{1}$。若 $t_{2}$ \tkanswer{$=$} $t_{1}$，可将两小车的碰撞视为弹性碰撞。
\item
吸能材料性能测试。
将吸能材料紧贴于小车 2 的前端。重复步骤③。测得小车 2 通过光电门 A 的时间为 $10.00 \UmsT$，两车碰撞后，依次测得小车 1 和 2 通过光电门 B 的时间分别为 $15.00 \UmsT$、$30.00 \UmsT$，不计吸能材料的质量，计算可得碰撞后两小车总动能与碰撞前小车 2 动能的比值为 \tkanswer{0.56}（结果保留 2 位有效数字）。
\end{steps}
\end{enumerate}

\twopicture[labelA=2025广东11a, labelB=2025广东11b, wA=3.50cm, wB=8.50cm, align=right]{figs/11a.png}{figs/11b.png}

%% answer:
\jdanswer{
\begin{enumerate}
\item
8.260
\item
②时间相等；③$=$；④$0.56$
\end{enumerate}
}
%% memo:
\memoanswer{
（1）根据题意，由图可知，小球的直径为 $d = 8\Umm + 26.0 \times 0.01\Umm = 8.260\Umm$

（2）②若已平衡小车在轨道甲上所受摩擦力及其他阻力，小车将在轨道甲上做匀速直线运动，通过两个光电门的速度相等，即通过光电门 A 和 B 的时间相等。

③若两个小车发生弹性碰撞，由于两个小车的质量相等，则碰撞后两个小车的速度互换，即碰撞后小车 1 的速度等于碰撞前小车 2 的速度，则有 $t_{2} = t_{1}$

④根据题意可知，碰撞前小车 2 的速度为

$v_{0}=\frac{d}{\Delta t_{1}}=\frac{1\times10^{-2}}{10\times10^{-3}}\Ums=1\Ums$

碰撞后，小车 1 和小车 2 的速度分别为

$v_{1}=\frac{d}{\Delta t_{2}}=\frac{2}{3}\Ums$，$v_{2}=\frac{d}{\Delta t_{3}}=\frac{1}{3}\Ums$

则碰撞后两小车总动能与碰撞前小车 2 动能的比值为

$\frac{E_{\mathrm{k}}'}{E_{\mathrm{k}}}=\frac{\frac{1}{2}mv_{1}^{2}+\frac{1}{2}mv_{2}^{2}}{\frac{1}{2}mv_{0}^{2}}=\frac{5}{9}\approx0.56$
}

\item
%% number: 12
%% paperName: 2025年广东高考第 12 题 8 分
%% typeId: 4
%% type: 实验题
%% chapter: 电磁感应
%% point: 自感与互感，涡流
%% method: 无
%% score: 8
%% degree: 550
%% duplicateId: 0
%% body:
科技小组制作的涡流制动演示装置由电磁铁和圆盘控制部分组成。

图~\ref{2025广东12a} 是电磁铁磁感应强度的测量电路。所用器材有：电源 $E$（电动势 $15 \UV$，内阻不计）；电流表 A（量程有 $0.6 \UA$ 和 $3 \UA$，内阻不计）；滑动变阻器 $R_{\mathrm{P}}$（最大阻值 $100 \UO$）；定值电阻 $R_{0}$（阻值 $10 \UO$）；开关 S；磁传感器和测试仪；电磁铁（线圈电阻 $16 \UO$）；导线若干。图~\ref{2025广东12b} 是实物图，图中电机和底座相固定，圆形铝盘和电机转轴相固定。

请完成下列实验操作和计算。

\begin{enumerate}
\item
量程选择和电路连接。
\begin{steps}
\item
由器材参数可得电路中的最大电流为 \tkanswer{0.58} $\UA$（结果保留 2 位有效数字），为减小测量误差，电流表的量程选择 $0.6 \UA$ 挡。
\item
图~\ref{2025广东12b} 中已正确连接了部分电路，请在虚线框中完成 $R_{\mathrm{P}}$、$R_{0}$ 和 A 间的实物图连线。
\end{steps}
\item
磁感应强度 $B$ 和电流 $I$ 关系测量。
\begin{steps}
\item
将图~\ref{2025广东12a} 中的磁传感器置于电磁铁中心，滑动变阻器 $R_{\mathrm{P}}$ 的滑片 P 置于 b 端。置于 b 端目的是使电路中的电流 \tkanswer{最小}，保护电路安全。
\item
将滑片 P 缓慢滑到某一位置，闭合 S。此时 A 的示数如图~\ref{2025广东12c} 所示，读数为 \tkanswer{0.48} $\UA$。分别记录测试仪示数 $B$ 和 $I$，断开 S。
\item
保持磁传感器位置不变，重复步骤②。
\item
如图~\ref{2025广东12d} 是根据部分实验数据描绘的 $B-I$ 图线，其斜率为 \tkanswer{30} $\UmT/\UA$（结果保留 2 位有效数字）。
\end{steps}
\item
制动时间 $t$ 测量。

利用图~\ref{2025广东12b} 所示装置测量了 $t$，结果表明 $B$ 越大，$t$ 越小。
\end{enumerate}

\fourpicture[labelA=2025广东12a, labelB=2025广东12b, labelC=2025广东12c, labelD=2025广东12d, wA=5.50cm, wB=7.00cm, wC=3.50cm, wD=3.80cm, align=right]
{figs/12a.png}
{figs/12b.png}
{figs/12c.png}
{figs/12d.png}

%% answer:
\jdanswer{
\begin{enumerate}
\item
①$0.58$；②连线如图所示
\item
①最小；②$0.48$；③$30$
\item
—
\end{enumerate}
}
%% memo:
\memoanswer{
（1）①由题知，电源内阻不计、电流表内阻不计，则当滑动变阻器的阻值为零时，电路中有最大电流 $I = \frac{E}{R_{0} + R_{\text{线}}} = 0.58 \UA$。

②由于电路中最大电流为 $0.58 \UA$，则电流表应选择 $0 \sim 0.6 \UA$ 量程，根据电路图实物图连线如下：

\onepicture[w=5.00cm]{figs/12_an.png}

（2）①滑动变阻器 $R_{\mathrm{P}}$ 的滑片 P 置于 b 端时滑动变阻器的电阻最大，电路中的电流最小，保护电路安全。

②电流表读数为 $0.48 \UA$。

③根据题图中数据可知 $B-I$ 图线斜率为 $k = \frac{15.2 - 5.6}{0.5 - 0.18} \UmT/\UA = 30 \UmT/\UA$。

（3）利用图（b）所示装置测量了 $t$，结果表明 $B$ 越大，$t$ 越小。
}

\item
%% number: 13
%% paperName: 2025年广东高考第 13 题 9 分
%% typeId: 6
%% type: 计算题
%% chapter: 气体性质
%% point: 玻意耳定律
%% method: 无
%% score: 9
%% degree: 550
%% duplicateId: 0
%% body:
如图是某铸造原理示意图，往气室注入空气增加压强，使金属液沿升液管进入已预热的铸型室，待铸型室内金属液冷却凝固后获得铸件。柱状铸型室通过排气孔与大气相通，大气压强 $p_{0} = 1.0\times10^{5} \UPa$，铸型室底面积 $S_{1} = 0.2 \Umq$，高度 $h_{1} = 0.2 \Um$，底面与注气前气室内金属液面高度差 $H = 0.15 \Um$，柱状气室底面积 $S_{2} = 0.8 \Umq$，注气前气室内气体压强为 $p_{0}$，金属液的密度 $\rho = 5.0\times10^{3} \Ukgmc$，重力加速度取 $g = 10 \Umsq$，空气可视为理想气体，不计升液管的体积。

\begin{enumerate}
\item
求金属液刚好充满铸型室时，气室内金属液面下降的高度 $h_{2}$ 和气室内气体压强 $p_{1}$。
\item
若在注气前关闭排气孔使铸型室密封，且注气过程中铸型室内温度不变，求注气后铸型室内的金属液高度为 $h_{3} = 0.04 \Um$ 时，气室内气体压强 $p_{2}$。
\end{enumerate}

\onepicture[w=4.00cm, align=right]{figs/13.png}

%% answer:
\jdanswer{
\begin{enumerate}
\item
$h_{2} = 0.05 \Um$，$p_{1} = 1.2\times10^{5} \UPa$
\item
$p_{2} = 1.35\times10^{5} \UPa$
\end{enumerate}
}
%% memo:
\memoanswer{
（1）根据体积关系 $S_{1}h_{1} = S_{2}h_{2}$

可得下方液面下降高度 $h_{2} = 0.05 \Um$

此时下方气体的压强 $p_{1} = p_{0} + \rho g(h_{1} + H + h_{2})$

代入数据可得 $p_{1} = 1.2\times10^{5} \UPa$

（2）初始时，上方铸型室气体的压强为 $p_{0}$，体积 $V = S_{1}h_{1}$

当上方铸型室液面高为 $h_{3} = 0.04 \Um$ 时体积为 $V' = S_{1}(h_{1} - h_{3})$

根据玻意耳定律 $p_{0}V = p'V'$

可得此时上方铸型室液面高为 $h_{3} = 0.04 \Um$ 时气体的压强为 $p' = 1.25 \times 10^{5} \UPa$

同理根据体积关系 $S_{1}h_{3} = S_{2}h_{4}$

可得 $h_{4} = 0.01 \Um$

此时下方气室内气体压强 $p_{2} = p' + \rho g(H + h_{3} + h_{4})$

代入数据可得 $p_{2} = 1.35\times10^{5} \UPa$
}

\item
%% number: 14
%% paperName: 2025年广东高考第 14 题 13 分
%% typeId: 6
%% type: 计算题
%% chapter: 功和能
%% point: 功率
%% method: 无
%% score: 13
%% degree: 450
%% duplicateId: 0
%% body:
如图所示，用开瓶器取出紧塞在瓶口的软木塞时，先将拔塞钻旋入木塞内，随后下压把手，使齿轮绕固定支架上的转轴转动，通过齿轮啮合，带动与木塞相固定的拔塞钻向上运动。从 0 时刻开始，顶部与瓶口齐平的木塞从静止开始向上做匀加速直线运动，木塞所受摩擦力 $f$ 随位移大小 $x$ 的变化关系为 $f = f_{0}(1 - \frac{x}{h})$，其中 $f_{0}$ 为常量，$h$ 为圆柱形木塞的高，木塞质量为 $m$，底面积为 $S$，加速度为 $a$，齿轮半径为 $r$，重力加速度为 $g$，瓶外气压减瓶内气压为 $\Delta p$ 且近似不变，瓶子始终静止在桌面上。（提示：可用 $f-x$ 图线下的“面积”表示 $f$ 所做的功）求：

\begin{enumerate}
\item
木塞离开瓶口的瞬间，齿轮的角速度 $\omega$。
\item
拔塞的全过程，拔塞钻对木塞做的功 $W$。
\item
拔塞过程中，拔塞钻对木塞作用力的瞬时功率 $P$ 随时间 $t$ 变化的表达式。
\end{enumerate}

\onepicture[w=9.00cm, align=right]{figs/14.png}

%% answer:
\jdanswer{
\begin{enumerate}
\item
$\omega = \frac{\sqrt{2ah}}{r}$
\item
$W = mah + mgh + \frac{1}{2}f_{0}h + \Delta pSh$
\item
$P = magt + ma^{2}t + \Delta pSat + f_{0}at - \frac{f_{0}a^{3}t^{3}}{2h}$
\end{enumerate}
}
%% memo:
\memoanswer{
（1）木塞的末速度等于齿轮线速度，对木塞，根据运动学公式 $v^{2} = 2ah$

根据角速度和线速度的关系 $v = \omega r$

联立可得 $\omega = \frac{\sqrt{2ah}}{r}$

（2）根据题意画出木塞摩擦力与运动距离的关系图如图所示。

\onepicture[w=3.50cm]{\includesvg{figs/14_an}}

可得摩擦力对木塞所做的功为 $W_{f} = - \frac{1}{2}f_{0}h$

对木塞，根据动能定理

$W + W_{f} - mgh - \Delta pSh = \frac{1}{2}mv^{2} - 0$

解得 $W = mah + mgh + \frac{1}{2}f_{0}h + \Delta pSh$

（3）设开瓶器对木塞的作用力为 $F$，对木塞，根据牛顿第二定律 $F - mg - f - \Delta pS = ma$

速度 $v = at$

位移 $x = \frac{1}{2}at^{2}$

开瓶器的功率 $P = Fv$

联立可得 $P = magt + ma^{2}t + \Delta pSat + f_{0}at - \frac{f_{0}a^{3}t^{3}}{2h}$
}

\item
%% number: 15
%% paperName: 2025年广东高考第 15 题 16 分
%% typeId: 6
%% type: 计算题
%% chapter: 电场
%% point: 带电粒子的直线运动
%% method: 分情况讨论
%% score: 16
%% degree: 350
%% duplicateId: 0
%% body:
如图是研究颗粒碰撞荷电特性装置的简化图。两块水平绝缘平板与两块竖直的平行金属平板相接。金属平板之间接高压电源产生匀强电场。一带电颗粒从上方绝缘平板左端 A 点处，由静止开始向右下方运动，与下方绝缘平板在 B 点处碰撞，碰撞时电荷量改变，反弹后离开下方绝缘平板瞬间，颗粒的速度与所受合力垂直，其水平分速度与碰前瞬间相同，竖直分速度大小变为碰前瞬间的 $k$ 倍（$k < 1$）。已知颗粒质量为 $m$，两绝缘平板间的距离为 $h$，两金属平板间的距离为 $d$，B 点与左平板的距离为 $l$，电源电压为 $U$，重力加速度为 $g$。忽略空气阻力和电场的边缘效应。求：

\begin{enumerate}
\item
颗粒碰撞前的电荷量 $q$。
\item
颗粒在 B 点碰撞后的电荷量 $Q$。
\item
颗粒从 A 点开始运动到第二次碰撞过程中，电场力对它做的功 $W$。
\end{enumerate}

\onepicture[w=3.80cm, align=right]{figs/15.png}

%% answer:
\jdanswer{
\begin{enumerate}
\item
$q = \frac{mgdl}{Uh}$
\item
$Q = \frac{kmgdh}{Ul}$
\item
若 $(4k + 1)l + \frac{4k^{3}h^{2}}{l} \leq d$，$W = \frac{mgl^{2}}{h} + 4k^{2}mgh + \frac{4k^{4}mgh^{3}}{l^{2}}$；

若 $(4k + 1)l + \frac{4k^{3}h^{2}}{l} > d$，$W = \frac{mgl^{2}}{h} + \frac{kmgh(d - l)}{l}$
\end{enumerate}
}
%% memo:
\memoanswer{
（1）根据题意可知，颗粒在竖直方向上做自由落体，则有 $h = \frac{1}{2}gt^{2}$

水平方向上做匀加速直线运动，则有 $\frac{qU}{d} = ma$，$l = \frac{1}{2}at^{2}$

解得 $q = \frac{mgdl}{Uh}$

（2）根据题意可知，颗粒与绝缘板第一次碰撞时，竖直分速度为 $v_{y1} = \sqrt{2gh}$

水平分速度为 $v_{x1} = \sqrt{2al} = \sqrt{\frac{2gl^{2}}{h}}$

则第一次碰撞后竖直分速度为 $v_{y2} = kv_{y1} = k\sqrt{2gh}$

设第一次碰撞后颗粒速度方向与水平方向夹角为 $\theta$，则有

$\tan\theta = \frac{v_{y2}}{v_{x1}} = \frac{kh}{l}$

由于第一次碰撞后瞬间颗粒所受合力与速度方向垂直，则有

$\tan\theta = \frac{\frac{UQ}{d}}{mg} = \frac{UQ}{mgd}$

联立解得 $Q = \frac{kmgdh}{Ul}$

（3）根据题意可知，由于 $k < 1$，则第一次碰撞后颗粒不能返回上绝缘板，若颗粒第二次碰撞是和下绝缘板碰撞，设从第一次碰撞后到第二次碰撞前的运动时间为 $t'$，则有

$t' = \frac{2k\sqrt{2gh}}{g} = 2k\sqrt{\frac{2h}{g}}$

水平方向上做匀加速直线运动，加速度为 $a' = \frac{UQ}{md} = \frac{kgh}{l}$

水平方向运动的距离为 $l' = v_{x1}t' + \frac{1}{2}a't'^{2} = 4kl + \frac{4k^{3}h^{2}}{l}$

则电场对颗粒做的功为 $W = \frac{qUl}{d} + \frac{QUl'}{d} = \frac{mgl^{2}}{h} + 4k^{2}mgh + \frac{4k^{4}mgh^{3}}{l^{2}}$

若 $l + l' > d$，则颗粒第二次碰撞是和右侧金属板碰撞，则颗粒从第一次碰撞到第二次碰撞过程中水平方向位移为 $d - l$，颗粒从 A 点开始运动到第二次碰撞过程中，电场对颗粒做的功为 $W = \frac{qUl}{d} + \frac{QU(d - l)}{d} = \frac{mgl^{2}}{h} + \frac{kmgh(d - l)}{l}$
}

\end{enumerate}

\end{document}
